We ran twenty models under the protocol of Section~\ref{sec:method}.
This section describes each family in turn; the numerical results are
reported in Section~\ref{sec:results}.

\subsection{Naive Baselines}

We need a floor below which any forecasting model should not fall, and
that floor is the naive family. We ran five baselines:
\textbf{historical mean}, \textbf{last value}, \textbf{drift},
\textbf{seasonal-naive} (at $s=7$), and \textbf{ARIMA naive-seasonal}
(an ARIMA(3,1,3) applied to the de-seasonalised series, used as a
weak baseline in some related work). The point of the baselines is to
quantify how much of the variance in the series can be explained
without any real learning; any forecasting model that does not
comfortably beat them is not worth running.

\subsection{Classical (S)ARIMA / ARIMAX Family}

We ran an exhaustive AIC grid search over $(p, d, q)(P, D, Q)_7$ with
$p, q \in [0, 4]$, $P, Q \in [0, 2]$, $d = D = 1$ (so the seasonal
difference at lag 7 is applied) and $s = 7$. The grid was repeated
with and without exogenous variables (temperature, weekend flag,
holiday flag), giving two families: SARIMA (no exog) and SARIMAX (with
exog). Models are fit via maximum likelihood with the
\texttt{statsmodels} SARIMAX backend, with rolling-origin refits at
each test window so the model has access to the most recent actuals.

In addition to the SARIMA grid, we evaluated the simpler members of
the Box-Jenkins family that the course rubric explicitly requires:
\textbf{ARMA}, \textbf{ARIMA} (two different orders),
\textbf{AutoARIMA} (seasonal and non-seasonal, via \texttt{pmdarima}),
and \textbf{ARIMAX} (ARIMA with exogenous variables). The reason for
including ARIMAX as a separate model rather than as a SARIMAX variant
is that it does not include an explicit seasonal term: it lets the
exogenous variables (especially the weekend / holiday flags) absorb
most of the weekly cycle. As we will see in Section~\ref{sec:results},
this turns out to be a competitive design on Boulder data.

\subsection{LSTM}

The LSTM (\texttt{scripts/lstm\_train.py}) uses five input features
(energy, mean temperature, weekend flag, holiday flag, month) and
performs a hyper-parameter sweep over sequence length
$\in \{14, 28\}$, hidden size $\in \{32, 64\}$, and number of layers
$\in \{1, 2\}$. Each configuration is trained with Adam, MSE loss,
dropout 0.1 and early stopping on the 60-day validation split. The
winning configuration is \texttt{seq14\_h64\_l2}: 14-day sequences,
hidden size 64, 2 layers. At inference time the LSTM forecasts one
day ahead and is rolled forward autoregressively for the full 7-day
horizon, feeding back its own prediction with the known future
exogenous row each step.

\subsection{Zero-Shot Foundation Models}

We benchmark four foundation models in zero-shot mode, no fine-tuning,
just feed the history and read the forecast:
\begin{itemize}[leftmargin=*]
    \item \textbf{Google TimesFM 1.0} (200M parameters) and
    \textbf{TimesFM 2.0} (500M)~\cite{das2024timesfm}, both
    decoder-only transformers with patch-based input embedding,
    pre-trained on a public time-series corpus.
    \item \textbf{Amazon Chronos-T5 base} (200M) and
    \textbf{Chronos-T5 large} (710M)~\cite{ansari2024chronos}, T5
    encoder-decoder models that quantise the input into tokens and
    sample probabilistic forecasts. We take the median of 20 samples
    as the point forecast.
\end{itemize}
These four models are evaluated under exactly the same rolling-origin
protocol as everything else. Comparing them lets us test two
sub-questions in parallel: (a) Google or Amazon, on this kind of
series? (b) Does bigger help inside each family?

\subsection{PatchTST From Scratch}

For the SoTA push we implemented PatchTST~\cite{nie2023patchtst} from
scratch in PyTorch (Fig.~\ref{fig:patchtst}). The architecture
choices, made deliberately small to match our $\sim$1900-point
dataset:
\begin{itemize}[leftmargin=*]
    \item Context length: 56 days (8 weeks).
    \item Patch length: 7 days, non-overlapping stride 7 $\Rightarrow$
    8 tokens per input.
    \item 3 transformer encoder layers, $d_{\text{model}}=96$, 4
    attention heads, GELU activation, dropout 0.15.
    \item \textbf{Channel-mixed patching}: each patch is a
    $(7, 4)$ block over the four input channels (energy + 3 exog),
    flattened to 28-dim and linearly embedded; cheaper and more
    data-efficient than channel-independent attention for our
    sequence length.
    \item \textbf{RevIN}~\cite{kim2022revin} on the target channel
    only: per-sample mean / std subtraction before the encoder and
    inverse transform on the output. This is what kept the model
    competitive against the regime shifts in the data (COVID drop,
    holiday weeks).
\end{itemize}
We train each model with AdamW (LR $5 \times 10^{-4}$, weight decay
$10^{-4}$), cosine annealing, early stopping with patience 25, batch
size 32. Training a single PatchTST takes $\approx 2$~minutes on the
RTX 5070~Ti.

\begin{figure}[!htbp]
\centering
\includegraphics[width=0.95\textwidth]{fig_patchtst}
\caption{PatchTST forward pass as implemented in
\texttt{scripts/patchtst\_train.py}. The energy channel goes through
RevIN before patching; the three exogenous channels are concatenated
into each patch and flattened together. The head is a single linear
layer that directly predicts the next 7 days, with no autoregression.}
\label{fig:patchtst}
\end{figure}

A single PatchTST trained on this setup overfits noticeably (train
MSE around $0.08$, validation MSE around $0.19$ in the standardised
domain) and the rolling-origin MAE comes out worse than ARIMAX.
The fix is the seed ensemble: we train three independent PatchTST
models (same architecture and hyper-parameters, different random seeds:
42, 7 and 2024) and average their predictions at inference. The
combined model is what we call the \textbf{PatchTST 3-seed ensemble}
in the leaderboard.

\subsection{Ensembles}

Finally, we test two ways of combining the base learners.

\noindent
\textbf{Simple mean (PatchTST + LSTM).} The two strongest deep
learning models in the project, picked \emph{a priori} by architecture
rather than by test score, averaged with equal weights:
\begin{equation*}
\text{Ensemble}_{\text{mean}} = \tfrac{1}{2} \big(\hat{y}_{\text{PatchTST}} + \hat{y}_{\text{LSTM}}\big).
\end{equation*}
This is the simplest defensible combination because it does not use
any test-set information to choose the weights.

\noindent
\textbf{Stacked meta-learners (Ridge and NNLS, LOWO).} As a check on
whether a more sophisticated combiner would help, we ran two leave-one-window-out
stacking variants: a \texttt{RidgeCV} with $\alpha \in
\{0.01, 0.1, 1, 10, 100\}$ and a non-negative convex combination
solved by SLSQP with weights summing to one. Both are trained on the
predictions of the six top base learners (SARIMA, SARIMAX, ARIMAX,
ARIMA, LSTM, PatchTST) across three windows and tested on the held-out
fourth. The result --- foreshadowing what we describe in
Section~\ref{sec:results} --- is that both meta-learners overfit
badly with only 21 training points per fold and end up worse than the
unweighted mean.
